\section{Úvod do statistiky}\label{sec:ai-uvod-statistika}

Data sama o~sobě nejsou hotovou informací. Číslo \(45\,000\) například získá význam až ve chvíli, kdy víme, že jde o~měsíční mzdu v~korunách, známe období a~skupinu lidí, ke které se vztahuje, a~rozumíme způsobu jeho výpočtu. Statistika poskytuje jazyk a~metody, jimiž lze větší množství dat uspořádat, shrnout a~porovnat. Ve strojovém učení je tato role zásadní: model se učí právě z~dat a~jejich nevhodná interpretace se proto promítne i~do jeho výsledků.

\begin{definition}\label{def:ai-statisticky-soubor}
    \emph{Statistický znak} je sledovaná vlastnost objektů. Statistický soubor s~\(n\) objekty a~\(p\) sledovanými znaky budeme zapisovat jako matici
    \[
        \mathbf{X}
        =
        \begin{pmatrix}
            x_{11} & x_{12} & \cdots & x_{1p}\\
            x_{21} & x_{22} & \cdots & x_{2p}\\
            \vdots & \vdots & \ddots & \vdots\\
            x_{n1} & x_{n2} & \cdots & x_{np}
        \end{pmatrix}.
    \]
    Hodnota \(x_{ij}\) je hodnota \(j\)-tého znaku u~\(i\)-tého objektu.
\end{definition}

Řádky matice tedy představují objekty a~sloupce jednotlivé statistické znaky. Pokud sledujeme pouze jeden znak, zapisujeme statistický soubor jako vektor
\[
    \mathbf{X}=(x_1,x_2,\ldots,x_n).
\]
Tučné písmo je důležité: symbol \(\mathbf{X}\) označuje celý statistický soubor, zatímco \(x_i\) označuje jednu jeho hodnotu.

\begin{figure}[ht]
    \centering
    \begin{tikzpicture}[every node/.style={font=\normalsize}]
    \node[task,fill=blue!10,minimum width=25mm] (objects) at (0,0)
        {objekty\\\(\mathbf{x}_1,\ldots,\mathbf{x}_n\)};
    \node[task,fill=orange!15,minimum width=25mm] (matrix) at (4,0)
        {\(\mathbf{X}=(x_{ij})\)};
    \node[task,fill=green!12,minimum width=28mm] (features) at (8,0)
        {statistické znaky\\\(\mathbf{X}_1,\ldots,\mathbf{X}_p\)};
    \draw[diredge] (objects) -- node[above]{řádky} (matrix);
    \draw[diredge] (matrix) -- node[above]{sloupce} (features);
\end{tikzpicture}

    \caption{Dva rovnocenné pohledy na statistický soubor.}
    \label{fig:ai-statisticky-soubor}
\end{figure}

Pro stejně dlouhé statistické soubory
\[
    \mathbf{X}=(x_1,\ldots,x_n),
    \qquad
    \mathbf{Y}=(y_1,\ldots,y_n)
\]
a~čísla \(a,b\in\R\) budeme používat po složkách definované operace
\[
    a\mathbf{X}+b\mathbf{Y}
    =
    (ax_1+by_1,\ldots,ax_n+by_n).
\]
Tento zápis využijeme například při posouvání, změně měřítka a~standardizaci dat.

\warningbox{Souhrnná charakteristika nikdy nepopisuje každý objekt zvlášť. Průměrná mzda není mzdou „běžného“ člověka, nižší inflace neznamená pokles cen a~výsledek za celou skupinu nemusí platit pro každou její část. Při interpretaci je vždy nutné znát kontext dat.}
